\section{Two exact Gaussian sign transitions below outer order minus one}
\label{sec:negative}
The incoming Euler moment identities prove $g_{s,b}<0$ for every
$s\ge-1$ and $b>0$ \cite[Section 6]{IncomingSharp}. They also exhibit
opposite signs farther down the outer-order ladder and ask for the
resulting sign changes. We resolve the first two integer cases
completely in the inner order $b$.

\subsection{A finite identity at every negative integer outer order}
At $a=0$, \eqref{ker:eq:continuation} simplifies to
\[
 F_{0,b}(z)=\frac{z\Li_b(z)}{1-z}=\Li_0(z)\Li_b(z).
\]
The product rule and $\D\Li_s(z)=\Li_{s-1}(z)$ give the following
exact identity, with no relation search.

\begin{proposition}[Negative-order product formula]\label{neg:prop:product}
For every integer $m\ge0$ and every $b>0$,
\begin{equation}\label{neg:eq:product}
 F_{-m,b}(z)=\sum_{j=0}^m\binom mj
                  \Li_{-j}(z)\Li_{b-m+j}(z)
\end{equation}
on the principal slit plane. In particular the Gaussian value is a
finite combination of shifted Dirichlet beta and eta values.
\end{proposition}
\begin{proof}
In the disk $F_{-m,b}=\D^mF_{0,b}$, by termwise differentiation.
Apply the finite Leibniz rule to the displayed product, then continue
analytically. No convergence of the original series at $z=\ii$ is
asserted or needed.
\end{proof}

For all complex $s$, analytic continuation of the residue-class identity
gives
\begin{equation}\label{neg:eq:Li-beta}
 \Li_s(\ii)=-2^{-s}\eta(s)+\ii\beta(s).
\end{equation}
Using
$\Li_0(\ii)=(-1+\ii)/2$,
$\Li_{-1}(\ii)=-1/2$,
$\Li_{-2}(\ii)=-\ii/2$,
and $\Li_{-3}(\ii)=1$ in \eqref{neg:eq:product} gives
\begin{align}
 2g_{-2,b}={}&-\beta(b-2)-2\beta(b-1)
              -2^{2-b}\eta(b-2)+2^{-b}\eta(b),\label{neg:eq:m2}\\
 2g_{-3,b}={}&2\beta(b)-\beta(b-3)-3\beta(b-2)
              -2^{3-b}\eta(b-3)
              +3\,2^{1-b}\eta(b-1).\label{neg:eq:m3}
\end{align}
The shifted arguments in these formulas can be nonpositive; they have
their continued values. The next theorem is a global sign assertion
about these finite combinations.

\subsection{Rational kernels and one-crossing Mellin integrals}
\begin{lemma}[A Mellin sign-change principle]\label{neg:lem:mellin}
Suppose a real function $w$ is positive on $(0,\tau)$ and negative
on $(\tau,\infty)$, up to null sets, and
\[
 J(b)=\int_0^\infty t^{b-1}w(t)\dd t
\]
and its first logarithmic derivative integral converge locally uniformly
for $b>0$. Then $\tau^{-b}J(b)$ is strictly decreasing on $(0,\infty)$.
Consequently $J$ has at most one positive zero. If it has one, that
zero is simple and $J'$ is negative there.
\end{lemma}
\begin{proof}
Differentiation gives
\[
 \frac{d}{db}\{\tau^{-b}J(b)\}
 =\tau^{-b}\int_0^\infty
       t^{b-1}\log(t/\tau)w(t)\dd t<0.
\]
The integrand is negative on both sides of $\tau$. At a zero of
$J$, the same inequality gives $J'<0$.
\end{proof}

\begin{theorem}[Complete first negative-order sign transitions]
\label{neg:thm:transitions}
For each $m\in\{2,3\}$ there is exactly one $b_m>0$ such that
$g_{-m,b_m}=0$. This zero is simple, and
\begin{equation}\label{neg:eq:signs}
 g_{-m,b}>0\quad(0<b<b_m),\qquad
 g_{-m,b}<0\quad(b>b_m).
\end{equation}
Moreover $0<b_2<1<b_3$. The two Dirichlet combinations
\eqref{neg:eq:m2} and \eqref{neg:eq:m3} have exactly the same
respective zero and sign descriptions.
\end{theorem}
\begin{proof}
Differentiate the axis case of \eqref{ker:eq:continuation} under its
bounded rational integrand. If $Y_b=e^{-T_b}$ and
$T_b\sim\operatorname{Gamma}(b,1)$, then
\[
 g_{-m,b}=\E q_m(Y_b),\qquad
 q_m(y)=\Ima\left[
   \D^m\frac{z^2}{(1-z)(1-zy)}\right]_{z=\ii}.
\]
For each fixed $m$, this differentiation is justified uniformly on a
complex neighborhood of $\ii$, because $0\le y\le1$ and the
denominators stay away from zero. Direct differentiation yields
\begin{align}
 q_2(y)&=\frac{(1+y)(y^4+6y^2-3)}{2(1+y^2)^3},\label{neg:eq:q2}\\
 q_3(y)&=-\frac{(1+y)(y^4-22y^2+1)}{(1+y^2)^4}.
                    \label{neg:eq:q3}
\end{align}
Each rational function has exactly one zero in $(0,1)$, namely
\[
 r_2=\sqrt{-3+2\sqrt3},\qquad
 r_3=\sqrt{11-2\sqrt{30}}.
\]
It is negative below that zero and positive above it. After the change
$y=e^{-t}$, write
\begin{equation}\label{neg:eq:mellin}
 \Gamma(b)g_{-m,b}=\int_0^\infty
          t^{b-1}e^{-t}q_m(e^{-t})\dd t.
\end{equation}
The kernel $w_m(t)=e^{-t}q_m(e^{-t})$ is positive before
$\tau_m=-\log r_m$ and negative after it. It is bounded at zero
and decays exponentially at infinity, so all local parameter
differentiations are legitimate. \Cref{neg:lem:mellin} proves at most
one simple crossing, from positive to negative.

As $b\downarrow0$, $T_b\to0$ in probability because $\E T_b=b$.
As $b\to\infty$, $T_b\to\infty$ in probability, for example by
Chebyshev's inequality using its mean $b$ and variance $b$.
Since $q_m$ is bounded and continuous on $[0,1]$, these limits give
\begin{align*}
 g_{-2,0+}=q_2(1)=\tfrac12,&\qquad
 g_{-2,\infty}=q_2(0)=-\tfrac32,\\
 g_{-3,0+}=q_3(1)=\tfrac52,&\qquad
 g_{-3,\infty}=q_3(0)=-1.
\end{align*}
Continuity therefore proves existence. Finally the exact values
\[
 g_{-2,1}=-\frac34+\frac{\log2}{4}<0,
 \qquad g_{-3,1}=1+\frac\pi4>0
\]
locate the two roots relative to one. They follow either from
\eqref{neg:eq:m2}--\eqref{neg:eq:m3} or from differentiating
$F_{1,1}(z)=\tfrac12\log^2(1-z)$. These two isolated evaluations
were already displayed in \cite[Section 6]{IncomingSharp}; the new
conclusion is the complete one-crossing classification for all $b>0$.
\end{proof}

Independent numerical evaluation of the Dirichlet formulas and the
Mellin integrals gives
\[
 b_2=0.4238611967707581269539978159\ldots,
 \qquad b_3=2.4369081043976327447535806729\ldots.
\]
These decimals are diagnostics; uniqueness, simplicity, and the sign
classification are analytic results. No small-residual test is used
to decide whether a zero exists.

\paragraph{Why the next integer orders need new analysis.}
The rational kernel construction and finite identity
\eqref{neg:eq:product} work for every $m$. A one-crossing argument
does not automatically persist. For example,
\[
 q_5(y)=-\frac{(1+y)
 (y^8-236y^6+1446y^4-236y^2+1)}{(1+y^2)^6}
\]
has two zeros in $(0,1)$. To see this exactly, divide its degree-eight
numerator by $y^4$ and set $T=y^2+y^{-2}>2$. The resulting quadratic
is $T^2-236T+1444$, whose distinct roots $118\pm8\sqrt{195}$
both exceed $2$. The map $y\mapsto T$ decreases bijectively from
$(0,1)$ onto $(2,\infty)$, giving exactly two simple zeros.
More general Mellin variation bounds give
an upper bound by the number of kernel sign changes, but existence
of all the permitted crossings is a separate problem. This is a
concrete continuation of the negative-order question, not a proposed
blanket extension of Gaussian negativity.
